\documentclass[11pt,a4paper]{article}
\usepackage[utf8]{inputenc}
\usepackage[T1]{fontenc}
\usepackage{amsmath,amssymb,amsthm}
\usepackage[margin=2.7cm]{geometry}
\usepackage{hyperref}
\hypersetup{colorlinks=true,linkcolor=blue,citecolor=blue,urlcolor=blue}

\newtheorem{theorem}{Theorem}
\newtheorem{corollary}[theorem]{Corollary}
\newtheorem{remark}[theorem]{Remark}

\newcommand{\Gcal}{\mathcal{G}}
\newcommand{\dd}{\mathrm{d}}
\newcommand{\sig}{\sigma_0}

\title{A regular energy identity for the scalar sector of stabilized
two-brane models: an alternative proof of the stability criterion,
valid without the monotonic-warp assumption}

\author{Hicham Boufourou\\[2pt]
\small Independent researcher, Ingelmunster, Belgium\\
\small \texttt{hicham.boufourou@hotmail.com}}

\date{2026}

\begin{document}
\maketitle

\begin{abstract}
For a five-dimensional two-brane model with a bulk scalar stabilizing
the interbrane distance, we derive an energy identity for the scalar
perturbation sector which is manifestly non-negative in the bulk and
entirely regular: it contains no factor of $A'$, no inverse
d'Alembertian, and no eigenvalue-dependent boundary condition. Its
natural boundary condition reproduces the certified junction
condition, which fixes the normalization of the brane term rather than
leaving it to be chosen. The identity yields the tachyon-free
criterion $s_i\,\sig''(y_i)/\sig'(y_i)\ge 0$ at both branes, which
reproduces the necessary and sufficient criterion of Lesgourgues and
Sorbo in its domain of validity. The point of the present derivation
is that it does \emph{not} assume the warp factor to be monotonic:
Lesgourgues and Sorbo explicitly exclude the symmetric configuration
$U_+=U_-$, for which $a(y)$ is non-monotonic, and that configuration
is precisely the one realized by a $\mathbb{Z}_2$-symmetric
Goldberger--Wise background, which we treat as an application.
\end{abstract}

\section{Introduction and statement of what is new}
\label{sec:intro}

Two-brane models with a bulk scalar are stable against scalar
perturbations under conditions that are known. Lesgourgues and
Sorbo~\cite{LS} derived, without approximation and for arbitrary bulk
and brane potentials, a necessary and sufficient criterion for the
absence of tachyonic modes, in terms of two boundary quantities
$g_\pm$ built from the background. Their analysis proceeds by a
phase-space (shooting) argument on a function $F(m^2)$ whose zeros are
the Kaluza--Klein masses. Independently, the interval approach of
Carena, Lykken and Park~\cite{CLP} provides the action-principle
framework in which the boundary equations follow from a variation with
metric fluctuations that do not vanish on the branes, including the
correct treatment of the degenerate tensor decomposition at null
four-momentum. The energy-identity technique used below has a direct
precedent in the tensor sector, where Mukohyama and Kofman~\cite{MK}
establish positivity by exactly this method.

Against this background, two things in the present note are not
already in the literature, and nothing else is claimed.

\begin{enumerate}
\item \textbf{A constructive identity.} We exhibit an exact identity
(Theorem~\ref{thm:main}) in which the bulk contribution is a manifest
square and the boundary contribution collapses to a single term whose
sign is read off directly. The derivation is local, uses no integrating
factor, and is self-validating: the natural boundary condition of the
resulting functional is the certified junction condition, so the
normalization of the brane term is fixed by the calculation rather than
assumed.

\item \textbf{Validity without the monotonic-warp assumption.} The
theorem of~\cite{LS} assumes $a(y)$ monotonic ($H\neq 0$ throughout).
Its authors state explicitly that the case $U_+=U_-$ makes $a(y)$
non-monotonic between the branes and that their analysis does not hold
there. Our identity makes no assumption on the warp: the factor $A'$
cancels identically. Section~\ref{sec:app} treats the
$\mathbb{Z}_2$-symmetric Goldberger--Wise background, which falls
exactly in that gap.
\end{enumerate}

A methodological remark is recorded in Section~\ref{sec:disc}: a
projection of the boundary condition onto scalar variations degenerates
on the null cone. This is an error one can make; the interval
methodology of~\cite{CLP} avoids it by construction, since the boundary
equation is imposed as a tensor equation and only afterwards
decomposed.

\section{Setup}
\label{sec:setup}

We work on $M_4\times[0,L]$ with two physical boundaries. The action is
\begin{equation}
S=\int \dd^4x\,\dd y\,\sqrt{-g}\Big[2M_5^3R-\tfrac12(\partial\sigma)^2
-V(\sigma)\Big]
+4M_5^3\!\oint_{\partial M}\!K
-\sum_i\int \dd^4x\,\sqrt{-\gamma}\,U_i(\sigma),
\label{eq:action}
\end{equation}
with outward-normal orientations $s_0=-1$ at $y=0$ and $s_L=+1$ at
$y=L$. Boundary equations follow from the variation of
\eqref{eq:action} with fluctuations that need not vanish on the
boundaries, in the sense of~\cite{CLP}; we do not impose junction
conditions independently.

The background is
$\dd s^2=e^{2A(y)}\eta_{\mu\nu}\dd x^\mu \dd x^\nu+\dd y^2$,
$\sigma=\sig(y)$, with
\begin{equation}
A''=-\frac{\sig'^2}{3M_5^3},\qquad
6M_5^3A'^2=\tfrac12\sig'^2-V(\sig),\qquad
\sig''=V'(\sig)-4A'\sig' .
\label{eq:bg}
\end{equation}
Throughout we take the brane potentials to be linear in the scalar,
$U_i(\sigma)=T_i+J_i\sigma$, so that $\dd^2U_i/\dd\sigma^2=0$; the
background junctions are then
$U_i=3s_iM_5^3A'(y_i)$ and $J_i=-s_i\sig'(y_i)$.

\paragraph{Scalar perturbations.}
In the gauge in which the metric perturbation is conformal and
$g_{\mu5}$ is removed, the sector is described by two functions
$(f,s)$ --- the conformal factor and the scalar fluctuation --- subject
to the momentum constraint
\begin{equation}
\mathcal{C}\equiv 3M_5^3\big(f'+2A'f\big)+\sig'\,s=0,
\label{eq:C1}
\end{equation}
and to the boundary condition obtained by variation,
\begin{equation}
s'+2\sig'f\Big|_{y_i}=0 .
\label{eq:BC}
\end{equation}
The four-dimensional mass eigenvalue problem is
$\mathbb{A}X=m^2\mathbb{K}X$ on the constraint surface, with
$\mathbb{K}=e^{2A}\,\mathrm{diag}(3M_5^3,\tfrac12)$ positive definite,
so that
\begin{equation}
\|X\|^2_{\mathbb{K}}=\int_0^L e^{2A}\Big(3M_5^3f^2+\tfrac12 s^2\Big)\dd y>0 .
\label{eq:norm}
\end{equation}

\section{The identity}
\label{sec:identity}

Define the combination
\begin{equation}
\boxed{\;\Gcal:=\big(s'+2\sig'f\big)-\frac{\sig''}{\sig'}\,s\;}
\label{eq:G}
\end{equation}
which contains no $A'$ and requires only $\sig'\neq0$.

\begin{theorem}[Regular energy identity]
\label{thm:main}
On the constraint surface \eqref{eq:C1}, and with the brane terms of
\eqref{eq:action},
\begin{equation}
\boxed{\;
m^2\,\|X\|^2_{\mathbb{K}}
=\frac12\int_0^L e^{4A}\,\Gcal^2\,\dd y
+\sum_i s_i\left(\frac{e^{4A}\sig''}{2\sig'}\right)_{y_i} s(y_i)^2 \; }
\label{eq:main}
\end{equation}
\end{theorem}

\begin{proof}[Sketch; the symbolic verification is in the companion script]
Writing $V_{\rm bulk}$ for the potential density obtained from the
reduced Lagrangian, one verifies first that its Euler--Lagrange
derivatives reproduce the certified operator $\mathbb{A}$ in both
directions. Eliminating $f'$ and $f''$ by \eqref{eq:C1}, one then has
the exact algebraic identity
\begin{equation}
2V_{\rm bulk}=\tfrac12 e^{4A}\Gcal^2+\frac{\dd}{\dd y}\Big[\Gcal_y\Big],
\qquad
\Gcal_y=28M_5^3e^{4A}A'f^2-\tfrac43\sig'e^{4A}fs
+\frac{e^{4A}\sig''}{2\sig'}s^2,
\label{eq:tot}
\end{equation}
where the total derivative is taken \emph{using} \eqref{eq:C1}; the
residue vanishes identically, and independently on an explicit exact
background. The brane contribution, restricted to the constraint
surface, is $s_ie^{4A}\big(14M_5^3A'f^2-\tfrac23\sig'fs\big)$, and
\begin{equation}
s_i\Gcal_y-2\,\mathrm{GB}\big|_{\rm img}
=s_i\left(\frac{e^{4A}\sig''}{2\sig'}\right)s(y_i)^2,
\end{equation}
all dependence on $f$ cancelling identically by
$\sig''=V'(\sig)-4A'\sig'$. Collecting gives \eqref{eq:main}.
\end{proof}

\begin{remark}[Self-validation]
Varying the right-hand side of \eqref{eq:main} with respect to $s$ at
the boundary gives $e^{4A}\Gcal+2\varepsilon_is=0$ with
$\varepsilon_i=e^{4A}\sig''/(2\sig')$, which reduces \emph{exactly} to
$s'+2\sig'f=0$, i.e.\ to \eqref{eq:BC}. The normalization of the brane
term is therefore fixed and checked, not chosen.
\end{remark}

\begin{corollary}[Criterion]
\label{cor:crit}
Since the bulk term in \eqref{eq:main} is non-negative and
$\|X\|^2_{\mathbb{K}}>0$,
\begin{equation}
\boxed{\;s_i\,\frac{\sig''(y_i)}{\sig'(y_i)}\;\ge\;0
\ \text{ at both branes}\quad\Longrightarrow\quad m^2\ge0 .\;}
\end{equation}
With $s_0=-1$, $s_L=+1$: $\sig''(0)/\sig'(0)\le0$ and
$\sig''(L)/\sig'(L)\ge0$.
\end{corollary}

\begin{corollary}[Recovery of the published criterion]
For linear brane potentials, the boundary quantities of~\cite{LS} are
$g_i=-e^{A}\sig''(y_i)/\sig'(y_i)$ after conversion from conformal to
proper coordinates, so their condition $g_+<0$, $g_->0$ is the
orientation-covariant statement $s_ig_i<0$, identical to
Corollary~\ref{cor:crit}. Their third condition $\varphi_0'\neq0$ is
our chart condition $\sig'\neq0$. If in addition $\sig$ is monotonic,
the criterion reduces to $\sig''(0)\sig''(L)\le0$.
\end{corollary}

\section{Extension: non-monotonic warp}
\label{sec:app}

Equation~\eqref{eq:main} contains no $A'$. In particular it holds when
$A'$ vanishes at an interior point, a case excluded from the theorem
of~\cite{LS}. Two structural facts make this the interesting case
rather than a curiosity.

\begin{remark}
$A'\equiv0$ is impossible for a non-trivial stabilizer: by
\eqref{eq:bg}, $A'\equiv0$ forces $\sig'\equiv0$. Since $A''<0$
strictly, $A'$ is strictly decreasing and vanishes at most once on
$[0,L]$.
\end{remark}

\paragraph{$\mathbb{Z}_2$-symmetric Goldberger--Wise background.}
Take $\sig(y)=(q/m_\sigma)\sinh[m_\sigma(y-L/2)]/\cosh(x/2)$,
$x=m_\sigma L$. Then $\sig$ is odd about $y=L/2$, $\sig'$ is even, the
right-hand side of the constraint in \eqref{eq:bg} is even, hence
$A'^2$ is even; combined with strict monotonicity of $A'$ this forces
$A'$ odd about $L/2$, so $A'(L/2)=0$ and $A'(0)=-A'(L)>0$. The warp is
therefore non-monotonic, and the analysis of~\cite{LS} does not apply.
Our criterion does: $\sig''/\sig'=m_\sigma\tanh[m_\sigma(y-L/2)]$, so
\begin{equation}
s_0\frac{\sig''}{\sig'}\Big|_{0}=s_L\frac{\sig''}{\sig'}\Big|_{L}
=+m_\sigma\tanh(x/2)\;\ge\;0,
\end{equation}
and the background is tachyon-free.

\section{Discussion}
\label{sec:disc}

\paragraph{What this note does and does not claim.}
The criterion of Corollary~\ref{cor:crit} is not new: it is the
criterion of~\cite{LS} specialized to linear brane potentials. The
interval framework, the action principle with non-vanishing boundary
fluctuations, and the degenerate decomposition at $p^2=0$ are those
of~\cite{CLP}. The energy-identity technique has its precedent
in~\cite{MK}. What is offered here is the constructive identity
\eqref{eq:main}, with its self-validating boundary term, and its
validity on non-monotonic warps.

\paragraph{A methodological caveat, recorded for its own sake.}
If one imposes the $(\mu\nu)$ boundary equation by projecting it onto
scalar field variations rather than as a tensor equation, the
determinant of the resulting conditions on the trace and
$\partial_\mu\partial_\nu$ parts is proportional to $p^4$. At null
four-momentum the projection therefore degenerates, and a genuine
boundary condition silently becomes an empty identity. The interval
treatment of~\cite{CLP} is immune to this, since there the boundary
equation is a tensor equation which is decomposed only afterwards --- see
their equations (3.53)--(3.55) and (4.43)--(4.44). We record the point
because the failure mode is not visible in the projected equations
themselves.

\paragraph{Open.}
The sub-branch $\sig''(y_i)=0$ makes the boundary term of
\eqref{eq:main} vanish identically, so $m^2\ge0$ trivially; but it is
also the locus where a massless mode exists~\cite{LS,MK}, and the
massless sector there requires separate treatment.

\begin{thebibliography}{9}
\bibitem{LS} J.~Lesgourgues and L.~Sorbo,
\emph{Goldberger--Wise variations: stabilizing brane models with a bulk
scalar}, Phys.\ Rev.\ D \textbf{69} (2004) 084010,
\href{https://arxiv.org/abs/hep-th/0310007}{arXiv:hep-th/0310007}.

\bibitem{CLP} M.~Carena, J.~Lykken and M.~Park,
\emph{The interval approach to braneworld gravity},
\href{https://arxiv.org/abs/hep-ph/0506305}{arXiv:hep-ph/0506305}.

\bibitem{MK} S.~Mukohyama and L.~Kofman,
\emph{Brane gravity at low energy}, Phys.\ Rev.\ D \textbf{65} (2002)
124025, \href{https://arxiv.org/abs/hep-th/0112115}{arXiv:hep-th/0112115}.

\bibitem{GW} W.~D.~Goldberger and M.~B.~Wise,
\emph{Modulus stabilization with bulk fields}, Phys.\ Rev.\ Lett.\
\textbf{83} (1999) 4922.

\bibitem{TM} T.~Tanaka and X.~Montes,
\emph{Gravity in the brane-world for two-branes model with stabilized
modulus}, Nucl.\ Phys.\ B \textbf{582} (2000) 259,
\href{https://arxiv.org/abs/hep-th/0001092}{arXiv:hep-th/0001092}.

\bibitem{multi} \emph{Stability of multibrane models} (2024),
\href{https://arxiv.org/abs/2408.15343}{arXiv:2408.15343}.
\end{thebibliography}

\end{document}
